\section{Propósito, alcance y cambios respecto de la 2.ª edición}
\label{sec:proposito}

\subsection{Para qué sirve este documento}

La pregunta de investigación del proyecto se formuló así:
\emph{técnicas de sintonía en entornos simulados como punto de partida en la
sintonía de drones reales}. Una sintonía hecha en simulación solo transfiere
al vehículo real si el simulador contiene, con números creíbles, los
fenómenos que limitan al controlador: el rezago de los motores, la saturación
de los actuadores, el acoplamiento entre inclinación y traslación, el ruido y
la vibración que ven los sensores, el error del estimador de actitud y el
viento. Este documento describe uno por uno esos fenómenos tal como están
programados en el paquete \cod{quadsim/}, explica por qué se modelaron así y
qué se dejó por fuera, y muestra la evidencia numérica de que el código hace
lo que dicen las ecuaciones.

La regla de redacción es simple. Cada ecuación de este texto corresponde a una
línea de código identificable (se cita el archivo y la función), y cada cifra
presentada como resultado sale del guion
\cod{experiments/model\_verification.py}, que la calcula, la imprime y la
guarda en \cod{docs/modelo/resultados\_verificacion.json}. Lo que es un
supuesto aparece en recuadros naranjas; lo que es un detalle de
implementación que conviene conocer antes de presentar el proyecto, en recuadros
azules.

\subsection{Alcance}

El simulador modela un cuadricóptero como un sólido rígido de seis grados de
libertad con cuatro rotores de dinámica propia (17 estados), en marco
inercial NED y marco cuerpo FRD, con:

\begin{itemize}
\item propulsor ESC + motor sin escobillas de primer orden, empuje y par
cuadráticos, zona muerta y asignación de control con desaturación por
prioridades;
\item aerodinámica del fuselaje (arrastre cuadrático y amortiguamiento
rotacional), viento como velocidad relativa del aire y efecto suelo opcional;
\item sensores con ruido de hoja de datos, sesgo, deriva, vibración a la
frecuencia de paso de pala, retardos y tasas de muestreo propias;
\item filtros biquad sobre el giroscopio, tres estimadores de actitud
(Madgwick, Mahony, EKF) y un filtro de Kalman de navegación;
\item el control de bajo nivel en cascada (ángulo $\to$ tasa $\to$ momento
normalizado) con ganancias calculadas del modelo por las reglas SIMC;
\item cinco controladores de posición/trayectoria (PID en cascada, MPC lineal,
MPC no lineal, aprendizaje por refuerzo puro y MPC con residuo aprendido) y
siete generadores de trayectoria;
\item perturbaciones: viento medio, turbulencia de Dryden, ráfaga $1-\cos$,
interferencia electromagnética en el magnetómetro, cambio de carga útil y
pérdida de eficiencia de un motor.
\end{itemize}

No modela la batería (el voltaje es constante), la elasticidad del bastidor,
la aerodinámica del rotor en vuelo de avance (\emph{blade flapping}, arrastre
inducido) ni la dinámica eléctrica del motor. La \cref{sec:limitaciones}
discute qué consecuencias tiene cada omisión.

\subsection{Qué cambió frente a la 2.ª edición}

La 2.ª edición del modelo (\cod{PAE\_Quadcopter\_GCS/simulator/physical\_model.py})
cumplió su papel didáctico, pero tenía cuatro limitaciones de fondo que
impedían usarla para el objetivo de investigación actual. La
\cref{tab:cambios} las resume.

\begin{table}[H]
\centering
\small
\caption{Diferencias de fondo entre la 2.ª y la 3.ª edición del modelo.}
\label{tab:cambios}
\begin{tabularx}{\textwidth}{>{\raggedright\arraybackslash}p{3.1cm}>{\raggedright\arraybackslash}X>{\raggedright\arraybackslash}X}
\toprule
Aspecto & 2.ª edición & 3.ª edición (este documento) \\
\midrule
Dinámica de rotación & Linealizada en hover y desacoplada por eje: $I\dot\omega=\tau-b_\omega\omega$, $\dot\theta=\omega$ por separado para roll, pitch y yaw & Newton--Euler completa con $\vect\omega\times J\vect\omega$, par giroscópico de los rotores y cuaternión (\cref{sec:dinamica}) \\
Traslación & Solo altura; inclinarse no desplazaba al dron y mantener una inclinación no costaba momento & Traslación 3D en NED con el empuje proyectado por la actitud completa \\
Integración & Euler semi-implícito a la tasa del lazo (500~Hz) & RK4 a \val{physHz}~Hz con el mando retenido entre muestras de control (\cref{sec:integracion}) \\
Marcos & FLU (z arriba), yaw como rumbo con signo invertido a mano & NED/FRD como PX4 y Quan Quan: yaw es rumbo de forma natural \\
Sensores y estimación & Ruido sumado al ángulo verdadero & IMU, magnetómetro, barómetro y posición externa con vibración; Madgwick, Mahony, EKF y Kalman de navegación \\
Vehículo por defecto & F450 de Jiřinec, 1.474~kg, en X & PAE-X450 (datos completos de CopterSim/RflySim, \cref{sec:parametros}) \\
\bottomrule
\end{tabularx}
\end{table}

\paragraph{Problemas del catálogo anterior.} La revisión del vehículo de
Jiřinec destapó dos errores que conviene declarar porque un evaluador atento
los puede encontrar. El primero es de masa. Jiřinec identifica su coeficiente
de empuje con la condición de hover $b=mg/(4\omega_0^2)$, con
$b=1.5108\times10^{-5}$~N/(rad/s)$^2$ y $\omega_0=463.1$~rad/s (su ec.~4.9).
Despejando la masa de esa misma ecuación se obtiene
$m = 4b\omega_0^2/g = \val{jirM}$~kg, no los 1.474~kg de la suma de
componentes de su tabla~4.2. Con 1.474~kg el hover exigiría
$\omega_0=\val{jirW}$~rad/s: los números de la tesis no son coherentes entre
sí, y la 2.ª edición usaba la combinación incoherente. El segundo error es de
geometría: las ecuaciones de momento de Jiřinec (sec.~2.2.4) están escritas
para bastidor en cruz (+), y la 2.ª edición lo cargaba en X. La 3.ª edición
corrige ambos (\cod{JIRINEC\_PLUS} en \cod{vehicles.py}, masa 1.321~kg, bastidor
+) y, sobre todo, cambia el vehículo de referencia por uno cuyos parámetros
forman un conjunto completo y verificable: el PAE-X450.

\paragraph{Lo que se conserva.} La cascada de control y su sintonía SIMC, la
convención de momento normalizado de PX4 en el lazo de tasa, el modelo de
actuador de primer orden y los filtros biquad pasan sin cambios conceptuales
de una edición a otra. Lo nuevo es la planta sobre la que actúan.

\paragraph{Correcciones del núcleo durante la verificación.} Al escribir
este documento, la verificación numérica encontró problemas en el propio
núcleo de la 3.ª edición que se corrigieron antes de cerrarlo (5 de octubre
de 2026): el acelerómetro llegaba sin filtrar al estimador de actitud y la
vibración se rectificaba en un sesgo de inclinación; el arranque en el aire
no inicializaba los retardos de los sensores, el estimador ni la covarianza
de navegación; y varios valores por defecto (derivativo del lazo de tasa,
filtro del giroscopio, ruidos del EKF, $\beta$ de Madgwick, integral del PID
de posición, estimador por defecto) se ajustaron con los barridos. El texto
describe el código corregido; las cifras históricas obtenidas con el núcleo
anterior se señalan como tales donde aparecen.

\paragraph{Revisión externa.} El 9 de octubre de 2026 una revisión técnica
independiente leyó el código y los resultados con instrucciones de buscar
errores. Sus hallazgos cambiaron varias partes de este documento. La
compensación de aceleración del AHRS resultó autorreferente cuando no hay
sensor de posición (\cref{sec:leishman}). Los retardos del ESC y del cómputo,
omitidos, pesan más de lo que la edición anterior estimaba
(\cref{sec:integracion}). La lectura del banco de controladores sobrevaloraba
lo que tres semillas permiten afirmar, y atribuía a la brecha de modelo un
error del aprendizaje por refuerzo que en su mayor parte ya existe en el
entorno de entrenamiento (\cref{sec:banco}). Lo que se pudo corregir en el
código se corrigió; lo que depende de medir el dron real quedó en la
\cref{sec:limitaciones}, con la lista de datos pendientes.

\begin{nota}[Cómo reproducir las cifras de este documento]
Desde la raíz del repositorio, \cod{bash docs/modelo/build.sh --verify}
corre la verificación numérica completa (bloques A a M, alrededor de un
minuto en un núcleo), regenera las figuras y \cod{valores.tex}, compila este
PDF y exporta \cod{docs/MODELO\_MATEMATICO.md}. Sin \cod{--verify} reutiliza
el último \cod{resultados\_verificacion.json}. Las cifras de los bancos de
controladores (\cref{sec:banco}) y de los barridos (\cref{sec:barridos})
vienen de sus propios guiones, citados en cada caso, porque requieren
cientos de corridas.
\end{nota}
